\chapter{Diamagnetism and Paramagnetism}

All materials respond to magnetic fields. Diamagnets are repelled (negative susceptibility); paramagnets are attracted (positive susceptibility). These responses, though typically weak, provide fundamental information about electronic structure.

\section{Langevin Diamagnetism}

All atoms have a diamagnetic response arising from the orbital motion of electrons in the applied field. The Langevin result for the molar diamagnetic susceptibility is:
\begin{equation}
  \chi_{\text{dia}} = -\frac{N_A \mu_0 Z e^2}{6m_e}\langle r^2 \rangle
\end{equation}

\begin{table}[H]
\centering
\caption{Molar diamagnetic susceptibilities ($10^{-6}$ \si{\centi\metre\cubed\per\mole}) of selected atoms and ions.}
\label{tab:diamag}
\begin{tabular}{@{}lSlS@{}}
\toprule
Species & {$\chi_{\text{dia}}$} & Species & {$\chi_{\text{dia}}$} \\
\midrule
He    & -1.9  & Cu$^+$  & -12 \\
Ne    & -7.2  & Ge      & -76 \\
Ar    & -19.3 & Si      & -13 \\
Kr    & -28.0 & NaCl    & -30 \\
Xe    & -43.0 & Diamond & -5.9 \\
\bottomrule
\end{tabular}
\end{table}

Diamagnetic susceptibilities increase with atomic number (larger $\langle r^2 \rangle$), as clearly seen in the noble gas series.


\section{Quantum Theory of Paramagnetism}

For atoms with permanent magnetic moments (unpaired electrons), the paramagnetic susceptibility follows the \textbf{Curie law}:
\begin{equation}
  \chi = \frac{C}{T} = \frac{N \mu_0 p^2 \mu_B^2}{3k_BT}, \qquad p = g\sqrt{J(J+1)}
\end{equation}
where $p$ is the effective number of Bohr magnetons.

\subsection{Rare Earth Ions}

The rare earth (4$f$) ions provide the most dramatic test of the quantum theory of paramagnetism, because the 4$f$ electrons are well shielded from the crystal environment and the free-ion values of $J$ apply.

\begin{table}[H]
\centering
\caption{Effective magnetic moments $p_{\text{eff}}$ of trivalent rare earth ions (in Bohr magnetons). Comparison of calculated values (Hund's rules) with experiment. Data from Guertin et al.\ (1973).}
\label{tab:rare_earth}
\begin{tabular}{@{}llSlSS@{}}
\toprule
Ion & Config. & {$J$} & $g$ & {$p_{\text{calc}}$} & {$p_{\text{exp}}$} \\
\midrule
Ce$^{3+}$ & $4f^1$  & 5/2 & 6/7  & 2.54 & 2.4 \\
Pr$^{3+}$ & $4f^2$  & 4   & 4/5  & 3.58 & 3.5 \\
Nd$^{3+}$ & $4f^3$  & 9/2 & 8/11 & 3.62 & 3.5 \\
Sm$^{3+}$ & $4f^5$  & 5/2 & 2/7  & 0.84 & 1.5 \\
Gd$^{3+}$ & $4f^7$  & 7/2 & 2    & 7.94 & 8.0 \\
Dy$^{3+}$ & $4f^9$  & 15/2 & 4/3 & 10.63 & 10.6 \\
Er$^{3+}$ & $4f^{11}$ & 15/2 & 6/5 & 9.59 & 9.5 \\
Yb$^{3+}$ & $4f^{13}$ & 7/2 & 8/7 & 4.54 & 4.5 \\
\bottomrule
\end{tabular}
\end{table}

The agreement between Hund's rules and experiment is remarkable for most rare earths. The notable exception is Sm$^{3+}$, where the calculated and experimental values disagree due to the small energy gap between the $J = 5/2$ ground state and the $J = 7/2$ excited state (Van Vleck paramagnetism).

\subsection{Iron Group Ions and Crystal Field Splitting}

For the 3$d$ (iron group) ions, the 3$d$ electrons are \textit{not} shielded from the crystal field and the orbital angular momentum is often ``quenched.'' Penney and Schlapp (1932)~\cite{penney1932crystal} provided the first theoretical treatment of crystal field effects on rare earth and iron group susceptibilities.

In octahedral coordination, the orbital contribution to the magnetic moment is reduced, and the effective moment is closer to the ``spin-only'' value $p = 2\sqrt{S(S+1)}$ rather than the free-ion value.

\section{Paramagnetic Susceptibility of Conduction Electrons}

The conduction electrons in a metal contribute a temperature-independent \textbf{Pauli paramagnetic susceptibility}:
\begin{equation}
  \chi_{\text{Pauli}} = \mu_0 \mu_B^2 D(E_F)
\end{equation}

This is much smaller than the Curie paramagnetism of localized moments because only electrons within $\sim k_BT$ of $E_F$ can flip their spin. For typical metals, $\chi_{\text{Pauli}} \sim 10^{-6}$, comparable in magnitude to the diamagnetic susceptibility.

\vspace{1cm}
\noindent\rule{\textwidth}{0.4pt}
\textit{We now turn to ferromagnetism and antiferromagnetism---the cooperative magnetic states where interactions between moments produce spontaneous order.}
